\documentclass[aspectratio=169,8pt]{beamer}
\usetheme{Madrid}
\usepackage[utf8]{inputenc}
\usepackage[T1]{fontenc}
\usepackage{amsmath,amssymb}
\usepackage{booktabs}
\usepackage{enumitem}
\setlist[enumerate]{label=\Alph*.,leftmargin=*,itemsep=1pt,topsep=2pt}
\setlist[itemize]{leftmargin=*,itemsep=1pt,topsep=2pt}
\setbeamerfont{block body}{size=\tiny}
\setbeamerfont{block title}{size=\scriptsize}
\setbeamerfont{frametitle}{size=\footnotesize}
\setbeamerfont{normal text}{size=\tiny}
\setbeamertemplate{navigation symbols}{}
\setbeamertemplate{itemize item}{\tiny$\bullet$}
\setbeamertemplate{itemize subitem}{\tiny$\circ$}

\title[GARP RAI]{GARP Risk \& AI --- 150 Memory \& Theory Questions}
\subtitle{Formulas, Names, Definitions, Theorems}
\author{Unofficial}
\date{\today}

\begin{document}
	
	\begin{frame}
		\titlepage
	\end{frame}
	
	% =================================================================
	\section{Formulas and Mathematical Definitions}
	% =================================================================
	
	\begin{frame}{Q1 --- Euclidean Distance Formula}
		\begin{block}{Question}
			In the K-means algorithm, Euclidean distance between two points $P$ and $Q$ in $m$ dimensions is expressed as:
			\begin{enumerate}
				\item $\sum_{j=1}^{m} |x_{jQ} - x_{jP}|$
				\item $\sqrt{\sum_{j=1}^{m} (x_{jQ} - x_{jP})^2}$
				\item $\max_j |x_{jQ} - x_{jP}|$
				\item $\sum_{j=1}^{m} (x_{jQ} - x_{jP})$
			\end{enumerate}
		\end{block}
		\begin{exampleblock}{Answer: B}\end{exampleblock}
		\begin{block}{Explanation}
			Euclidean distance is the square root of the sum of squared differences = L2-norm. Option A is Manhattan distance (L1). Option C is Chebyshev distance. Option D is signed sum (not a metric).
		\end{block}
	\end{frame}
	
	\begin{frame}{Q2 --- Manhattan Distance Formula}
		\begin{block}{Question}
			The Manhattan distance between two points $P$ and $Q$ in $m$ dimensions is also known as:
			\begin{enumerate}
				\item The L2-norm
				\item The L1-norm
				\item The L$\infty$-norm
				\item The cosine norm
			\end{enumerate}
		\end{block}
		\begin{exampleblock}{Answer: B}\end{exampleblock}
		\begin{block}{Explanation}
			Manhattan distance = L1-norm = sum of absolute differences. L2-norm is Euclidean; L$\infty$-norm is Chebyshev; cosine norm is a similarity metric.
		\end{block}
	\end{frame}
	
	\begin{frame}{Q3 --- Silhouette Score Formula}
		\begin{block}{Question}
			The silhouette score for a data point $i$ is defined as:
			\begin{enumerate}
				\item $s_i = a_i - b_i$
				\item $s_i = \dfrac{b_i - a_i}{\max(a_i, b_i)}$
				\item $s_i = \dfrac{a_i + b_i}{2}$
				\item $s_i = \dfrac{a_i \cdot b_i}{a_i + b_i}$
			\end{enumerate}
		\end{block}
		\begin{exampleblock}{Answer: B}\end{exampleblock}
		\begin{block}{Explanation}
			Silhouette score = $(b_i - a_i)/\max(a_i, b_i)$ where $a_i$ is mean intra-cluster distance and $b_i$ is mean distance to nearest other cluster. Range is $[-1, 1]$.
		\end{block}
	\end{frame}
	
	\begin{frame}{Q4 --- Silhouette Score Range}
		\begin{block}{Question}
			The silhouette score $s_i$ for any data point lies within which range?
			\begin{enumerate}
				\item $[0, 1]$
				\item $[-1, 1]$
				\item $[-1, 0]$
				\item $[0, \infty)$
			\end{enumerate}
		\end{block}
		\begin{exampleblock}{Answer: B}\end{exampleblock}
		\begin{block}{Explanation}
			Silhouette scores are bounded in $[-1, 1]$. $s \approx 1$ = well-clustered; $s \approx 0$ = boundary; $s < 0$ = possibly mis-clustered.
		\end{block}
	\end{frame}
	
	\begin{frame}{Q5 --- WCSS Definition}
		\begin{block}{Question}
			The Within-Cluster Sum of Squares (WCSS) for cluster $k$ is defined as:
			\begin{enumerate}
				\item $\sum_{j=1}^{n_k} d_j$
				\item $\sum_{j=1}^{n_k} d_j^2$
				\item $\sqrt{\sum_{j=1}^{n_k} d_j}$
				\item $\max_j d_j$
			\end{enumerate}
		\end{block}
		\begin{exampleblock}{Answer: B}\end{exampleblock}
		\begin{block}{Explanation}
			WCSS$_k = \sum_{j=1}^{n_k} d_j^2$ where $d_j$ is the distance from point $j$ to the cluster centroid. Squared distances ensure positive and negative deviations do not cancel.
		\end{block}
	\end{frame}
	
	\begin{frame}{Q6 --- Total WCSS Formula}
		\begin{block}{Question}
			The total WCSS for all $K$ clusters is given by:
			\begin{enumerate}
				\item $\sum_{k=1}^{K} \text{WCSS}_k$
				\item $\prod_{k=1}^{K} \text{WCSS}_k$
				\item $\max_k \text{WCSS}_k$
				\item $\frac{1}{K} \sum_{k=1}^{K} \text{WCSS}_k$
			\end{enumerate}
		\end{block}
		\begin{exampleblock}{Answer: A}\end{exampleblock}
		\begin{block}{Explanation}
			Total WCSS = sum of individual cluster WCSS values. This is the objective function minimised by K-means.
		\end{block}
	\end{frame}
	
	\begin{frame}{Q7 --- Bias-Variance Decomposition}
		\begin{block}{Question}
			The expected total prediction error of a model can be decomposed as:
			\begin{enumerate}
				\item Bias + Variance
				\item Bias$^2$ + Variance$^2$
				\item Bias$^2$ + Variance + Irreducible Error
				\item Bias + Variance + Irreducible Error
			\end{enumerate}
		\end{block}
		\begin{exampleblock}{Answer: C}\end{exampleblock}
		\begin{block}{Explanation}
			Total error = Bias$^2$ + Variance + Irreducible Error. The irreducible component represents noise that no model can eliminate.
		\end{block}
	\end{frame}
	
	\begin{frame}{Q8 --- Regularization Loss (Ridge)}
		\begin{block}{Question}
			The Ridge regression loss function is:
			\begin{enumerate}
				\item $\frac{1}{N}\sum (y_i - \hat{y}_i)^2 + \lambda \sum |\beta_j|$
				\item $\frac{1}{N}\sum (y_i - \hat{y}_i)^2 + \lambda \sum \beta_j^2$
				\item $\sum (y_i - \hat{y}_i)^2$
				\item $\lambda \sum |\beta_j|$
			\end{enumerate}
		\end{block}
		\begin{exampleblock}{Answer: B}\end{exampleblock}
		\begin{block}{Explanation}
			Ridge uses L2 penalty: $\lambda \sum \beta_j^2$. LASSO uses L1 penalty: $\lambda \sum |\beta_j|$. Options A and C are LASSO and OLS respectively.
		\end{block}
	\end{frame}
	
	\begin{frame}{Q9 --- Regularization Loss (LASSO)}
		\begin{block}{Question}
			The LASSO regression loss function is:
			\begin{enumerate}
				\item $\frac{1}{N}\sum (y_i - \hat{y}_i)^2 + \lambda \sum \beta_j^2$
				\item $\frac{1}{N}\sum (y_i - \hat{y}_i)^2 + \lambda \sum |\beta_j|$
				\item $\sum \beta_j^2$
				\item $\sum (y_i - \hat{y}_i)^2$
			\end{enumerate}
		\end{block}
		\begin{exampleblock}{Answer: B}\end{exampleblock}
		\begin{block}{Explanation}
			LASSO uses L1 penalty: $\lambda \sum |\beta_j|$. It can drive coefficients to exactly zero, performing feature selection.
		\end{block}
	\end{frame}
	
	\begin{frame}{Q10 --- Elastic Net Loss Function}
		\begin{block}{Question}
			The Elastic Net regression loss function combines:
			\begin{enumerate}
				\item L1 and L2 penalties
				\item Only L1 penalty
				\item Only L2 penalty
				\item Neither penalty
			\end{enumerate}
		\end{block}
		\begin{exampleblock}{Answer: A}\end{exampleblock}
		\begin{block}{Explanation}
			Elastic Net = L1 + L2 = $\lambda_1 \sum \beta_j^2 + \lambda_2 \sum |\beta_j|$. It inherits sparsity from L1 and grouping from L2.
		\end{block}
	\end{frame}
	
	\begin{frame}{Q11 --- Precision Formula}
		\begin{block}{Question}
			Precision is defined as:
			\begin{enumerate}
				\item $\dfrac{TP}{TP + FN}$
				\item $\dfrac{TP}{TP + FP}$
				\item $\dfrac{TN}{TN + FP}$
				\item $\dfrac{TP + TN}{TP + TN + FP + FN}$
			\end{enumerate}
		\end{block}
		\begin{exampleblock}{Answer: B}\end{exampleblock}
		\begin{block}{Explanation}
			Precision = TP/(TP+FP). Option A is recall (sensitivity); Option C is specificity; Option D is accuracy.
		\end{block}
	\end{frame}
	
	\begin{frame}{Q12 --- Recall Formula}
		\begin{block}{Question}
			Recall (sensitivity) is defined as:
			\begin{enumerate}
				\item $\dfrac{TP}{TP + FP}$
				\item $\dfrac{TP}{TP + FN}$
				\item $\dfrac{TN}{TN + FP}$
				\item $\dfrac{FP}{FP + TN}$
			\end{enumerate}
		\end{block}
		\begin{exampleblock}{Answer: B}\end{exampleblock}
		\begin{block}{Explanation}
			Recall = TP/(TP+FN) = true positive rate. Option A = precision; Option C = specificity; Option D = false positive rate.
		\end{block}
	\end{frame}
	
	\begin{frame}{Q13 --- F1 Score Formula}
		\begin{block}{Question}
			The F1 score is defined as:
			\begin{enumerate}
				\item $\dfrac{P + R}{2}$
				\item $\dfrac{2 \cdot P \cdot R}{P + R}$
				\item $\sqrt{P \cdot R}$
				\item $P \cdot R$
			\end{enumerate}
		\end{block}
		\begin{exampleblock}{Answer: B}\end{exampleblock}
		\begin{block}{Explanation}
			F1 = harmonic mean of precision and recall = $2PR/(P+R)$. Option A is arithmetic mean; Option C is geometric mean; Option D is product.
		\end{block}
	\end{frame}
	
	\begin{frame}{Q14 --- Specificity Formula}
		\begin{block}{Question}
			Specificity is defined as:
			\begin{enumerate}
				\item $\dfrac{TP}{TP + FN}$
				\item $\dfrac{TN}{TN + FP}$
				\item $\dfrac{TP}{TP + FP}$
				\item $\dfrac{FP}{FP + TN}$
			\end{enumerate}
		\end{block}
		\begin{exampleblock}{Answer: B}\end{exampleblock}
		\begin{block}{Explanation}
			Specificity = TN/(TN+FP) = true negative rate. Options A, C, D are recall, precision, and FPR respectively.
		\end{block}
	\end{frame}
	
	\begin{frame}{Q15 --- Accuracy Formula}
		\begin{block}{Question}
			Accuracy is defined as:
			\begin{enumerate}
				\item $\dfrac{TP + TN}{TP + TN + FP + FN}$
				\item $\dfrac{TP}{TP + FP}$
				\item $\dfrac{TP}{TP + FN}$
				\item $\dfrac{FP + FN}{TP + TN}$
			\end{enumerate}
		\end{block}
		\begin{exampleblock}{Answer: A}\end{exampleblock}
		\begin{block}{Explanation}
			Accuracy = (TP+TN)/Total. Option D is the error rate.
		\end{block}
	\end{frame}
	
	\begin{frame}{Q16 --- Error Rate Formula}
		\begin{block}{Question}
			Error rate (1 - accuracy) can be expressed as:
			\begin{enumerate}
				\item $\dfrac{TP + TN}{Total}$
				\item $\dfrac{FP + FN}{Total}$
				\item $\dfrac{TP + FP}{Total}$
				\item $\dfrac{FN + TN}{Total}$
			\end{enumerate}
		\end{block}
		\begin{exampleblock}{Answer: B}\end{exampleblock}
		\begin{block}{Explanation}
			Error rate = (FP+FN)/Total = 1 $-$ accuracy. Options A, C, D are not error rates.
		\end{block}
	\end{frame}
	
	\begin{frame}{Q17 --- MSE Formula}
		\begin{block}{Question}
			Mean Squared Error is:
			\begin{enumerate}
				\item $\frac{1}{N}\sum (y_i - \hat{y}_i)$
				\item $\frac{1}{N}\sum |y_i - \hat{y}_i|$
				\item $\frac{1}{N}\sum (y_i - \hat{y}_i)^2$
				\item $\sqrt{\frac{1}{N}\sum (y_i - \hat{y}_i)^2}$
			\end{enumerate}
		\end{block}
		\begin{exampleblock}{Answer: C}\end{exampleblock}
		\begin{block}{Explanation}
			MSE = average squared error. Option A would sum to zero; Option B is MAE; Option D is RMSE.
		\end{block}
	\end{frame}
	
	\begin{frame}{Q18 --- MAE Formula}
		\begin{block}{Question}
			Mean Absolute Error is:
			\begin{enumerate}
				\item $\frac{1}{N}\sum (y_i - \hat{y}_i)^2$
				\item $\frac{1}{N}\sum |y_i - \hat{y}_i|$
				\item $\sqrt{\frac{1}{N}\sum (y_i - \hat{y}_i)^2}$
				\item $\frac{1}{N}\sum (y_i - \hat{y}_i)$
			\end{enumerate}
		\end{block}
		\begin{exampleblock}{Answer: B}\end{exampleblock}
		\begin{block}{Explanation}
			MAE = average absolute error. It is less sensitive to outliers than MSE.
		\end{block}
	\end{frame}
	
	\begin{frame}{Q19 --- RMSE Formula}
		\begin{block}{Question}
			Root Mean Square Error is:
			\begin{enumerate}
				\item $\frac{1}{N}\sum (y_i - \hat{y}_i)^2$
				\item $\frac{1}{N}\sum |y_i - \hat{y}_i|$
				\item $\sqrt{\frac{1}{N}\sum (y_i - \hat{y}_i)^2}$
				\item $\frac{1}{N}\sum (y_i - \hat{y}_i)$
			\end{enumerate}
		\end{block}
		\begin{exampleblock}{Answer: C}\end{exampleblock}
		\begin{block}{Explanation}
			RMSE = square root of MSE. It is expressed in the same units as the target variable.
		\end{block}
	\end{frame}
	
	\begin{frame}{Q20 --- MAPE Formula}
		\begin{block}{Question}
			Mean Absolute Percentage Error is:
			\begin{enumerate}
				\item $\frac{1}{N}\sum |y_i - \hat{y}_i|$
				\item $100 \cdot \frac{1}{N}\sum \left|\dfrac{y_i - \hat{y}_i}{y_i}\right|$
				\item $\frac{1}{N}\sum (y_i - \hat{y}_i)^2$
				\item $\sqrt{\frac{1}{N}\sum (y_i - \hat{y}_i)^2}$
			\end{enumerate}
		\end{block}
		\begin{exampleblock}{Answer: B}\end{exampleblock}
		\begin{block}{Explanation}
			MAPE = 100 $\times$ mean of $|(y_i - \hat{y}_i)/y_i|$. It is scale-independent but fails when actuals are near zero.
		\end{block}
	\end{frame}
	
	\begin{frame}{Q21 --- Simple Linear Regression Equation}
		\begin{block}{Question}
			The simple linear regression model is:
			\begin{enumerate}
				\item $y_i = \beta_0 \cdot \beta_1 x_i + u_i$
				\item $y_i = \beta_0 + \beta_1 x_i + u_i$
				\item $y_i = \beta_0^{\beta_1} + u_i$
				\item $y_i = \beta_0 / (\beta_1 x_i) + u_i$
			\end{enumerate}
		\end{block}
		\begin{exampleblock}{Answer: B}\end{exampleblock}
		\begin{block}{Explanation}
			Simple linear regression: $y_i = \beta_0 + \beta_1 x_i + u_i$. $\beta_0$ is the intercept, $\beta_1$ is the slope, $u_i$ is the error term.
		\end{block}
	\end{frame}
	
	\begin{frame}{Q22 --- Logistic Function}
		\begin{block}{Question}
			The logistic (sigmoid) function is defined as:
			\begin{enumerate}
				\item $f(z) = z$
				\item $f(z) = e^z$
				\item $f(z) = \dfrac{1}{1 + e^{-z}}$
				\item $f(z) = \ln(z)$
			\end{enumerate}
		\end{block}
		\begin{exampleblock}{Answer: C}\end{exampleblock}
		\begin{block}{Explanation}
			Sigmoid = $1/(1+e^{-z})$. Output is bounded in $(0, 1)$, making it suitable for probability estimation in logistic regression.
		\end{block}
	\end{frame}
	
	\begin{frame}{Q23 --- Logistic Regression Probability}
		\begin{block}{Question}
			In a logistic regression with $m$ features, the probability that $y_i = 1$ is:
			\begin{enumerate}
				\item $P_i = \beta_0 + \beta_1 x_{1i} + \cdots + \beta_m x_{mi}$
				\item $P_i = \dfrac{1}{1 + e^{-(\beta_0 + \beta_1 x_{1i} + \cdots + \beta_m x_{mi})}}$
				\item $P_i = e^{\beta_0 + \beta_1 x_{1i} + \cdots + \beta_m x_{mi}}$
				\item $P_i = \ln(\beta_0 + \beta_1 x_{1i} + \cdots + \beta_m x_{mi})$
			\end{enumerate}
		\end{block}
		\begin{exampleblock}{Answer: B}\end{exampleblock}
		\begin{block}{Explanation}
			Logistic regression applies the sigmoid function to the linear predictor: $P_i = 1/(1 + e^{-(\beta_0 + \beta X)})$.
		\end{block}
	\end{frame}
	
	\begin{frame}{Q24 --- Log-Likelihood for Logit}
		\begin{block}{Question}
			The log-likelihood function for the logit model is:
			\begin{enumerate}
				\item $\sum_i [y_i \ln F(y_i) + (1 - y_i) \ln(1 - F(y_i))]$
				\item $\sum_i y_i$
				\item $\sum_i (y_i - F(y_i))^2$
				\item $\prod_i F(y_i)$
			\end{enumerate}
		\end{block}
		\begin{exampleblock}{Answer: A}\end{exampleblock}
		\begin{block}{Explanation}
			Log-likelihood sums log-probabilities of observed outcomes. Multiplying probabilities turns into summing logs. Option D is the likelihood (before log); Option C is a squared-error form (not likelihood).
		\end{block}
	\end{frame}
	
	\begin{frame}{Q25 --- OLS Slope Formula}
		\begin{block}{Question}
			The OLS estimate of the slope in simple linear regression is:
			\begin{enumerate}
				\item $\hat{\beta}_1 = \bar{y} - \hat{\beta}_0 \bar{x}$
				\item $\hat{\beta}_1 = \dfrac{\sum y_i x_i - N \bar{y} \bar{x}}{\sum x_i^2 - N \bar{x}^2}$
				\item $\hat{\beta}_1 = \dfrac{\bar{y}}{\bar{x}}$
				\item $\hat{\beta}_1 = \sum (y_i - \bar{y})(x_i - \bar{x})$
			\end{enumerate}
		\end{block}
		\begin{exampleblock}{Answer: B}\end{exampleblock}
		\begin{block}{Explanation}
			OLS slope = covariance of $x$ and $y$ divided by variance of $x$. Option A is the intercept formula (rearranged). Option C is a naive ratio; Option D is missing the denominator.
		\end{block}
	\end{frame}
	
	\begin{frame}{Q26 --- OLS Intercept Formula}
		\begin{block}{Question}
			Given the OLS slope estimate $\hat{\beta}_1$, the intercept is computed as:
			\begin{enumerate}
				\item $\hat{\beta}_0 = \bar{y} - \hat{\beta}_1 \bar{x}$
				\item $\hat{\beta}_0 = \bar{y} + \hat{\beta}_1 \bar{x}$
				\item $\hat{\beta}_0 = \bar{y} / \hat{\beta}_1 \bar{x}$
				\item $\hat{\beta}_0 = \hat{\beta}_1 \bar{x}$
			\end{enumerate}
		\end{block}
		\begin{exampleblock}{Answer: A}\end{exampleblock}
		\begin{block}{Explanation}
			The regression line passes through the point of means: $\bar{y} = \hat{\beta}_0 + \hat{\beta}_1 \bar{x}$, so $\hat{\beta}_0 = \bar{y} - \hat{\beta}_1 \bar{x}$.
		\end{block}
	\end{frame}
	
	\begin{frame}{Q27 --- Residual Sum of Squares}
		\begin{block}{Question}
			The residual sum of squares (RSS) is:
			\begin{enumerate}
				\item $\sum (y_i - \hat{y}_i)$
				\item $\sum (y_i - \hat{y}_i)^2$
				\item $\sum |y_i - \hat{y}_i|$
				\item $\sqrt{\sum (y_i - \hat{y}_i)^2}$
			\end{enumerate}
		\end{block}
		\begin{exampleblock}{Answer: B}\end{exampleblock}
		\begin{block}{Explanation}
			RSS = sum of squared residuals. OLS minimises RSS. Option A would be zero for OLS by construction.
		\end{block}
	\end{frame}
	
	\begin{frame}{Q28 --- Standardization Formula}
		\begin{block}{Question}
			The z-score standardization formula is:
			\begin{enumerate}
				\item $z = \dfrac{x - \mu}{\sigma}$
				\item $z = \dfrac{x - \min}{\max - \min}$
				\item $z = \ln(x)$
				\item $z = x^2$
			\end{enumerate}
		\end{block}
		\begin{exampleblock}{Answer: A}\end{exampleblock}
		\begin{block}{Explanation}
			Standardization = $(x-\mu)/\sigma$. Option B is min-max normalization.
		\end{block}
	\end{frame}
	
	\begin{frame}{Q29 --- Min-Max Normalization Formula}
		\begin{block}{Question}
			The min-max normalization formula is:
			\begin{enumerate}
				\item $x' = \dfrac{x - \mu}{\sigma}$
				\item $x' = \dfrac{x - x_{\min}}{x_{\max} - x_{\min}}$
				\item $x' = \ln(x)$
				\item $x' = \dfrac{1}{x}$
			\end{enumerate}
		\end{block}
		\begin{exampleblock}{Answer: B}\end{exampleblock}
		\begin{block}{Explanation}
			Min-max scales values to $[0,1]$: $x' = (x - x_{\min})/(x_{\max} - x_{\min})$.
		\end{block}
	\end{frame}
	
	\begin{frame}{Q30 --- PCA Component Equation}
		\begin{block}{Question}
			The $j$-th principal component is expressed as:
			\begin{enumerate}
				\item $PC_j = b_{1j} P_1 + b_{2j} P_2 + \cdots + b_{mj} P_m$
				\item $PC_j = P_1 \cdot P_2 \cdots P_m$
				\item $PC_j = \max(P_1, P_2, \ldots, P_m)$
				\item $PC_j = \dfrac{P_1 + P_2 + \cdots + P_m}{m}$
			\end{enumerate}
		\end{block}
		\begin{exampleblock}{Answer: A}\end{exampleblock}
		\begin{block}{Explanation}
			A principal component is a linear combination of standardised predictors weighted by component loadings $b_{ij}$.
		\end{block}
	\end{frame}
	
	\begin{frame}{Q31 --- Cosine Similarity Formula}
		\begin{block}{Question}
			Cosine similarity between two vectors $P$ and $Q$ is:
			\begin{enumerate}
				\item $\|P - Q\|$
				\item $\dfrac{P \cdot Q}{\|P\| \|Q\|}$
				\item $\|P\| + \|Q\|$
				\item $\dfrac{P}{Q}$
			\end{enumerate}
		\end{block}
		\begin{exampleblock}{Answer: B}\end{exampleblock}
		\begin{block}{Explanation}
			Cosine similarity = dot product divided by the product of magnitudes. Range is $[-1, 1]$.
		\end{block}
	\end{frame}
	
	\begin{frame}{Q32 --- TF-IDF Formula}
		\begin{block}{Question}
			The TF-IDF weight of term $i$ in document $j$ is:
			\begin{enumerate}
				\item $\text{TF}_{ij} + \text{IDF}_i$
				\item $\text{TF}_{ij} \times \text{IDF}_i$
				\item $\text{TF}_{ij} / \text{IDF}_i$
				\item $\text{TF}_{ij} - \text{IDF}_i$
			\end{enumerate}
		\end{block}
		\begin{exampleblock}{Answer: B}\end{exampleblock}
		\begin{block}{Explanation}
			TF-IDF = term frequency $\times$ inverse document frequency. Term frequency captures within-document frequency; IDF captures rarity across the corpus.
		\end{block}
	\end{frame}
	
	\begin{frame}{Q33 --- IDF Definition}
		\begin{block}{Question}
			The inverse document frequency of term $i$ across $D$ documents is:
			\begin{enumerate}
				\item $\ln(D / d_i)$
				\item $D \cdot d_i$
				\item $D + d_i$
				\item $D - d_i$
			\end{enumerate}
		\end{block}
		\begin{exampleblock}{Answer: A}\end{exampleblock}
		\begin{block}{Explanation}
			IDF$_i = \ln(D/d_i)$ where $d_i$ is the number of documents containing term $i$. If term appears in all documents, IDF = 0.
		\end{block}
	\end{frame}
	
	\begin{frame}{Q34 --- L2 Vector Normalization}
		\begin{block}{Question}
			L2 normalization of a vector $v$ divides each element by:
			\begin{enumerate}
				\item $\sum v_i$
				\item $\sqrt{\sum v_i^2}$
				\item $\max(v_i)$
				\item $|v_1|$
			\end{enumerate}
		\end{block}
		\begin{exampleblock}{Answer: B}\end{exampleblock}
		\begin{block}{Explanation}
			L2 norm = $\sqrt{\sum v_i^2}$. Dividing by it produces a unit vector.
		\end{block}
	\end{frame}
	
	\begin{frame}{Q35 --- Bayes Rule}
		\begin{block}{Question}
			Bayes' rule is written as:
			\begin{enumerate}
				\item $P(c|d) = P(d|c) P(c) / P(d)$
				\item $P(c|d) = P(c) + P(d)$
				\item $P(c|d) = P(d|c) / P(c)$
				\item $P(c|d) = P(c) \cdot P(d)$
			\end{enumerate}
		\end{block}
		\begin{exampleblock}{Answer: A}\end{exampleblock}
		\begin{block}{Explanation}
			Bayes: posterior = likelihood $\times$ prior / evidence. Used in Naive Bayes classification.
		\end{block}
	\end{frame}
	
	\begin{frame}{Q36 --- Naive Bayes Classification Rule}
		\begin{block}{Question}
			The Naive Bayes classification rule (argmax form) is:
			\begin{enumerate}
				\item $\hat{c} = \arg\max_c P(c)$
				\item $\hat{c} = \arg\max_c P(w_1|c) P(w_2|c) \cdots P(w_N|c) P(c)$
				\item $\hat{c} = \arg\min_c P(c)$
				\item $\hat{c} = \sum_c P(c)$
			\end{enumerate}
		\end{block}
		\begin{exampleblock}{Answer: B}\end{exampleblock}
		\begin{block}{Explanation}
			Naive Bayes predicts the class maximising the product of conditional word probabilities and the class prior, exploiting the independence assumption.
		\end{block}
	\end{frame}
	
	\begin{frame}{Q37 --- Attention Score Formula}
		\begin{block}{Question}
			The scaled dot-product attention score between query $Q_i$ and key $K_j$ is:
			\begin{enumerate}
				\item $Q_i \cdot K_j$
				\item $\dfrac{Q_i \cdot K_j}{\sqrt{d_k}}$
				\item $\dfrac{Q_i + K_j}{2}$
				\item $|Q_i - K_j|$
			\end{enumerate}
		\end{block}
		\begin{exampleblock}{Answer: B}\end{exampleblock}
		\begin{block}{Explanation}
			Scaled dot-product attention divides the dot product by $\sqrt{d_k}$ where $d_k$ is key dimension. This prevents softmax from saturating.
		\end{block}
	\end{frame}
	
	\begin{frame}{Q38 --- Softmax Function}
		\begin{block}{Question}
			The softmax function is defined as:
			\begin{enumerate}
				\item $\text{softmax}(x_i) = \dfrac{e^{x_i}}{\sum_j e^{x_j}}$
				\item $\text{softmax}(x_i) = x_i / \sum_j x_j$
				\item $\text{softmax}(x_i) = e^{x_i}$
				\item $\text{softmax}(x_i) = \ln(x_i)$
			\end{enumerate}
		\end{block}
		\begin{exampleblock}{Answer: A}\end{exampleblock}
		\begin{block}{Explanation}
			Softmax normalises exponentials into a probability distribution. Used in attention and multiclass outputs.
		\end{block}
	\end{frame}
	
	\begin{frame}{Q39 --- Gradient Descent Update Rule}
		\begin{block}{Question}
			The gradient descent weight update rule is:
			\begin{enumerate}
				\item $w_i^{\text{new}} = w_i^{\text{old}} + \eta \dfrac{\partial L}{\partial w_i}$
				\item $w_i^{\text{new}} = w_i^{\text{old}} - \eta \dfrac{\partial L}{\partial w_i}$
				\item $w_i^{\text{new}} = w_i^{\text{old}} \cdot \eta$
				\item $w_i^{\text{new}} = \eta$
			\end{enumerate}
		\end{block}
		\begin{exampleblock}{Answer: B}\end{exampleblock}
		\begin{block}{Explanation}
			Gradient descent subtracts the learning-rate-scaled gradient. Option A would be gradient ascent.
		\end{block}
	\end{frame}
	
	\begin{frame}{Q40 --- Dynamic Learning Rate (Exponential Decay)}
		\begin{block}{Question}
			Exponential decay of the learning rate is:
			\begin{enumerate}
				\item $\eta_t = \eta_0 + k t$
				\item $\eta_t = \eta_0 e^{-kt}$
				\item $\eta_t = \eta_0 \cdot k t$
				\item $\eta_t = \eta_0 / k$
			\end{enumerate}
		\end{block}
		\begin{exampleblock}{Answer: B}\end{exampleblock}
		\begin{block}{Explanation}
			Exponential decay: $\eta_t = \eta_0 e^{-kt}$, where $k$ controls decay and $t$ is the epoch.
		\end{block}
	\end{frame}
	
	\begin{frame}{Q41 --- Dynamic Learning Rate (Inverse Decay)}
		\begin{block}{Question}
			Inverse decay of the learning rate is:
			\begin{enumerate}
				\item $\eta_t = \eta_0 e^{-kt}$
				\item $\eta_t = \eta_0 / (1 + kt)$
				\item $\eta_t = \eta_0 (1 + kt)$
				\item $\eta_t = \eta_0 \cdot k$
			\end{enumerate}
		\end{block}
		\begin{exampleblock}{Answer: B}\end{exampleblock}
		\begin{block}{Explanation}
			Inverse decay: $\eta_t = \eta_0/(1+kt)$. Alternative to exponential decay for stochastic gradient descent.
		\end{block}
	\end{frame}
	
	% =================================================================
	\section{Names, Theorems, and Attribution}
	% =================================================================
	
	\begin{frame}{Q42 --- K-means Algorithm Attribution}
		\begin{block}{Question}
			The K-means algorithm is also known as:
			\begin{enumerate}
				\item Lloyd's algorithm
				\item Newton's algorithm
				\item Bayes' algorithm
				\item Fisher's algorithm
			\end{enumerate}
		\end{block}
		\begin{exampleblock}{Answer: A}\end{exampleblock}
		\begin{block}{Explanation}
			K-means is also known as Lloyd's algorithm, after Stuart Lloyd who developed it at Bell Labs in 1957.
		\end{block}
	\end{frame}
	
	\begin{frame}{Q43 --- Impossibility Theorem Attribution}
		\begin{block}{Question}
			The fairness impossibility theorem showing that demographic parity, predictive rate parity, and equal opportunity cannot all hold under unequal base rates is attributed to:
			\begin{enumerate}
				\item Kleinberg, Mullainathan, and Raghavan
				\item Fisher, Neyman, and Pearson
				\item Bentham and Mill
				\item Kant and Aristotle
			\end{enumerate}
		\end{block}
		\begin{exampleblock}{Answer: A}\end{exampleblock}
		\begin{block}{Explanation}
			Kleinberg, Mullainathan, and Raghavan (2016) proved this impossibility result. Chouldechova independently derived a related result.
		\end{block}
	\end{frame}
	
	\begin{frame}{Q44 --- Utilitarianism Founders}
		\begin{block}{Question}
			Which pair of philosophers are considered the founders of utilitarianism?
			\begin{enumerate}
				\item Kant and Aristotle
				\item Bentham and Mill
				\item Rawls and Nozick
				\item Plato and Socrates
			\end{enumerate}
		\end{block}
		\begin{exampleblock}{Answer: B}\end{exampleblock}
		\begin{block}{Explanation}
			Jeremy Bentham and John Stuart Mill developed classical utilitarianism in the 18th and 19th centuries.
		\end{block}
	\end{frame}
	
	\begin{frame}{Q45 --- Categorical Imperative Attribution}
		\begin{block}{Question}
			The concept of the categorical imperative is attributed to:
			\begin{enumerate}
				\item Immanuel Kant
				\item John Stuart Mill
				\item Jeremy Bentham
				\item Aristotle
			\end{enumerate}
		\end{block}
		\begin{exampleblock}{Answer: A}\end{exampleblock}
		\begin{block}{Explanation}
			Kant's categorical imperative: act only on maxims that could be universal laws.
		\end{block}
	\end{frame}
	
	\begin{frame}{Q46 --- Virtue Ethics Origin}
		\begin{block}{Question}
			Virtue ethics is most closely associated with which ancient philosopher?
			\begin{enumerate}
				\item Aristotle
				\item Kant
				\item Mill
				\item Bentham
			\end{enumerate}
		\end{block}
		\begin{exampleblock}{Answer: A}\end{exampleblock}
		\begin{block}{Explanation}
			Virtue ethics has roots in ancient Greek philosophy, particularly Aristotle's Nicomachean Ethics.
		\end{block}
	\end{frame}
	
	\begin{frame}{Q47 --- Transformer Paper Title}
		\begin{block}{Question}
			The 2017 paper introducing the transformer architecture is titled:
			\begin{enumerate}
				\item ``Attention Is All You Need''
				\item ``Deep Learning for AI''
				\item ``Neural Networks and Learning''
				\item ``Machine Learning Fundamentals''
			\end{enumerate}
		\end{block}
		\begin{exampleblock}{Answer: A}\end{exampleblock}
		\begin{block}{Explanation}
			``Attention Is All You Need'' by Vaswani et al. (2017) introduced the transformer architecture based on self-attention.
		\end{block}
	\end{frame}
	
	\begin{frame}{Q48 --- BERT Full Name}
		\begin{block}{Question}
			BERT stands for:
			\begin{enumerate}
				\item Bidirectional Encoder Representations from Transformers
				\item Bayesian Encoder for Recursive Tasks
				\item Binary Encoder Representation Technique
				\item Bidirectional Encoder for Random Text
			\end{enumerate}
		\end{block}
		\begin{exampleblock}{Answer: A}\end{exampleblock}
		\begin{block}{Explanation}
			BERT = Bidirectional Encoder Representations from Transformers. It is an encoder-only model released by Google in 2018.
		\end{block}
	\end{frame}
	
	\begin{frame}{Q49 --- GPT Full Name}
		\begin{block}{Question}
			GPT stands for:
			\begin{enumerate}
				\item General Purpose Transformer
				\item Generative Pre-trained Transformer
				\item Gradient Pre-training Technique
				\item Global Parameter Training
			\end{enumerate}
		\end{block}
		\begin{exampleblock}{Answer: B}\end{exampleblock}
		\begin{block}{Explanation}
			GPT = Generative Pre-trained Transformer. It is a decoder-only autoregressive model family from OpenAI.
		\end{block}
	\end{frame}
	
	\begin{frame}{Q50 --- BART Full Name}
		\begin{block}{Question}
			BART stands for:
			\begin{enumerate}
				\item Bidirectional and Auto-Regressive Transformer
				\item Bayesian Auto-Regressive Technique
				\item Binary Auto-Regressive Transformer
				\item Bidirectional Attention for Recognition Tasks
			\end{enumerate}
		\end{block}
		\begin{exampleblock}{Answer: A}\end{exampleblock}
		\begin{block}{Explanation}
			BART = Bidirectional and Auto-Regressive Transformer. It combines an encoder and decoder and is used for both understanding and generation tasks.
		\end{block}
	\end{frame}
	
	\begin{frame}{Q51 --- LSTM Full Name}
		\begin{block}{Question}
			LSTM stands for:
			\begin{enumerate}
				\item Long Short-Term Memory
				\item Linear Sequential Time Model
				\item Layered Stochastic Transformation Model
				\item Low-Signal Temporal Method
			\end{enumerate}
		\end{block}
		\begin{exampleblock}{Answer: A}\end{exampleblock}
		\begin{block}{Explanation}
			LSTM = Long Short-Term Memory. It is a type of RNN that mitigates the vanishing gradient problem via gating mechanisms.
		\end{block}
	\end{frame}
	
	\begin{frame}{Q52 --- DBSCAN Full Name}
		\begin{block}{Question}
			DBSCAN stands for:
			\begin{enumerate}
				\item Density-Based Spatial Clustering of Applications with Noise
				\item Data-Based Sample Clustering Algorithm for Networks
				\item Dynamic Bayesian Segmentation Clustering Algorithm
				\item Distributed Batch Spectral Clustering Approach
			\end{enumerate}
		\end{block}
		\begin{exampleblock}{Answer: A}\end{exampleblock}
		\begin{block}{Explanation}
			DBSCAN = Density-Based Spatial Clustering of Applications with Noise. It identifies core, border, and noise points using $\varepsilon$ and minPts.
		\end{block}
	\end{frame}
	
	\begin{frame}{Q53 --- LASSO Full Name}
		\begin{block}{Question}
			LASSO stands for:
			\begin{enumerate}
				\item Least Absolute Shrinkage and Selection Operator
				\item Linear Analysis of Statistical Sparse Outcomes
				\item Layered Adaptive Sparse Subset Optimization
				\item Least Squares Solution Operator
			\end{enumerate}
		\end{block}
		\begin{exampleblock}{Answer: A}\end{exampleblock}
		\begin{block}{Explanation}
			LASSO = Least Absolute Shrinkage and Selection Operator. It uses an L1 penalty and can zero out coefficients.
		\end{block}
	\end{frame}
	
	\begin{frame}{Q54 --- KNN Full Name}
		\begin{block}{Question}
			KNN stands for:
			\begin{enumerate}
				\item K-Nearest Neighbors
				\item Kernel Neural Network
				\item K-Normalized Nodes
				\item Knowledge Notation Network
			\end{enumerate}
		\end{block}
		\begin{exampleblock}{Answer: A}\end{exampleblock}
		\begin{block}{Explanation}
			KNN = K-Nearest Neighbors. A non-parametric classifier/regressor based on nearest training examples.
		\end{block}
	\end{frame}
	
	\begin{frame}{Q55 --- SVM Full Name}
		\begin{block}{Question}
			SVM stands for:
			\begin{enumerate}
				\item Support Vector Machine
				\item Supervised Vector Model
				\item Statistical Variance Method
				\item Sparse Vector Mechanism
			\end{enumerate}
		\end{block}
		\begin{exampleblock}{Answer: A}\end{exampleblock}
		\begin{block}{Explanation}
			SVM = Support Vector Machine. It finds the maximum-margin hyperplane separating classes.
		\end{block}
	\end{frame}
	
	\begin{frame}{Q56 --- CBoW Full Name}
		\begin{block}{Question}
			CBoW stands for:
			\begin{enumerate}
				\item Continuous Bag of Words
				\item Complex Bag of Words
				\item Categorical Bag of Words
				\item Continuous Batch of Words
			\end{enumerate}
		\end{block}
		\begin{exampleblock}{Answer: A}\end{exampleblock}
		\begin{block}{Explanation}
			CBoW = Continuous Bag of Words. One of two Word2Vec architectures: many inputs, one output.
		\end{block}
	\end{frame}
	
	\begin{frame}{Q57 --- RAG Full Name}
		\begin{block}{Question}
			RAG stands for:
			\begin{enumerate}
				\item Retrieval Augmented Generation
				\item Random Attention Gate
				\item Recurrent Adversarial Graph
				\item Regularized Attention Gradient
			\end{enumerate}
		\end{block}
		\begin{exampleblock}{Answer: A}\end{exampleblock}
		\begin{block}{Explanation}
			RAG = Retrieval Augmented Generation. It augments LLM responses with retrieved documents for factual accuracy.
		\end{block}
	\end{frame}
	
	\begin{frame}{Q58 --- RLHF Full Name}
		\begin{block}{Question}
			RLHF stands for:
			\begin{enumerate}
				\item Reinforcement Learning from Human Feedback
				\item Recursive Layered Hidden Features
				\item Random Loss with High Frequency
				\item Regularized Linear Hierarchical Framework
			\end{enumerate}
		\end{block}
		\begin{exampleblock}{Answer: A}\end{exampleblock}
		\begin{block}{Explanation}
			RLHF = Reinforcement Learning from Human Feedback. Used in LLM training to align outputs with human preferences.
		\end{block}
	\end{frame}
	
	\begin{frame}{Q59 --- PCA Full Name}
		\begin{block}{Question}
			PCA stands for:
			\begin{enumerate}
				\item Principal Component Analysis
				\item Predictive Cluster Aggregation
				\item Pattern Classification Algorithm
				\item Probabilistic Component Approximation
			\end{enumerate}
		\end{block}
		\begin{exampleblock}{Answer: A}\end{exampleblock}
		\begin{block}{Explanation}
			PCA = Principal Component Analysis. A dimensionality reduction technique using orthogonal linear combinations.
		\end{block}
	\end{frame}
	
	\begin{frame}{Q60 --- SR 11-7 Attribution}
		\begin{block}{Question}
			SR 11-7, ``Guidance on Model Risk Management'', was issued by which U.S. regulator?
			\begin{enumerate}
				\item The Federal Reserve and the OCC
				\item The SEC
				\item The FDIC alone
				\item The CFTC
			\end{enumerate}
		\end{block}
		\begin{exampleblock}{Answer: A}\end{exampleblock}
		\begin{block}{Explanation}
			SR 11-7 (2011) was issued jointly by the Federal Reserve and the Office of the Comptroller of the Currency (OCC) as supervisory guidance on model risk management.
		\end{block}
	\end{frame}
	
	\begin{frame}{Q61 --- GDPR Article on Breach Notification}
		\begin{block}{Question}
			Which GDPR article requires notification of a personal data breach to the supervisory authority within 72 hours?
			\begin{enumerate}
				\item Article 17
				\item Article 33
				\item Article 6
				\item Article 22
			\end{enumerate}
		\end{block}
		\begin{exampleblock}{Answer: B}\end{exampleblock}
		\begin{block}{Explanation}
			Article 33 = breach notification. Article 17 = right to erasure; Article 6 = lawful basis; Article 22 = automated decision-making rights.
		\end{block}
	\end{frame}
	
	\begin{frame}{Q62 --- GDPR Article on Right to Erasure}
		\begin{block}{Question}
			The GDPR right to erasure (right to be forgotten) is contained in:
			\begin{enumerate}
				\item Article 6
				\item Article 17
				\item Article 33
				\item Article 5
			\end{enumerate}
		\end{block}
		\begin{exampleblock}{Answer: B}\end{exampleblock}
		\begin{block}{Explanation}
			Article 17 = right to erasure. Article 5 = principles; Article 6 = lawful basis; Article 33 = breach notification.
		\end{block}
	\end{frame}
	
	\begin{frame}{Q63 --- NIST AI RMF Functions}
		\begin{block}{Question}
			The NIST AI Risk Management Framework is organised around four core functions:
			\begin{enumerate}
				\item Train, Test, Deploy, Monitor
				\item Govern, Map, Measure, Manage
				\item Design, Build, Audit, Retire
				\item Plan, Do, Check, Act
			\end{enumerate}
		\end{block}
		\begin{exampleblock}{Answer: B}\end{exampleblock}
		\begin{block}{Explanation}
			NIST AI RMF = Govern (structures), Map (context), Measure (assessment), Manage (mitigation).
		\end{block}
	\end{frame}
	
	\begin{frame}{Q64 --- Three Lines of Defence Attribution}
		\begin{block}{Question}
			The Three Lines of Defence model in financial risk management is best associated with:
			\begin{enumerate}
				\item The Institute of Internal Auditors (IIA)
				\item The Federal Reserve alone
				\item The Basel Committee alone
				\item The SEC alone
			\end{enumerate}
		\end{block}
		\begin{exampleblock}{Answer: A}\end{exampleblock}
		\begin{block}{Explanation}
			The IIA formalised the Three Lines of Defence model, later updated to the Three Lines Model in 2020.
		\end{block}
	\end{frame}
	
	% =================================================================
	\section{Definitions and Terminology}
	% =================================================================
	
	\begin{frame}{Q65 --- Inscrutability Definition}
		\begin{block}{Question}
			Inscrutability, as used in AI risk, is best defined as:
			\begin{enumerate}
				\item The inability to understand the model's internal mechanisms
				\item The inability to explain a single prediction
				\item The inability to generalise out-of-sample
				\item The inability to measure accuracy
			\end{enumerate}
		\end{block}
		\begin{exampleblock}{Answer: A}\end{exampleblock}
		\begin{block}{Explanation}
			Inscrutability = opaque internal mechanisms. Explainability = specific decisions. Generalisability = out-of-sample performance.
		\end{block}
	\end{frame}
	
	\begin{frame}{Q66 --- Interpretability Definition}
		\begin{block}{Question}
			Interpretability refers to:
			\begin{enumerate}
				\item Providing reasons for a specific decision
				\item Understanding the model's internal logic and mechanisms
				\item Measuring accuracy
				\item Reducing features
			\end{enumerate}
		\end{block}
		\begin{exampleblock}{Answer: B}\end{exampleblock}
		\begin{block}{Explanation}
			Interpretability = internal understanding. Explainability = specific-decision justification.
		\end{block}
	\end{frame}
	
	\begin{frame}{Q67 --- Explainability Definition}
		\begin{block}{Question}
			Explainability is best defined as:
			\begin{enumerate}
				\item Understanding internal model logic
				\item Providing human-understandable reasons for specific decisions
				\item Model training efficiency
				\item Feature reduction
			\end{enumerate}
		\end{block}
		\begin{exampleblock}{Answer: B}\end{exampleblock}
		\begin{block}{Explanation}
			Explainability = post-hoc justification of specific decisions in human terms.
		\end{block}
	\end{frame}
	
	\begin{frame}{Q68 --- Transparency Definition}
		\begin{block}{Question}
			Transparency in AI refers to:
			\begin{enumerate}
				\item Understanding internal model logic
				\item Openness about the model's design, algorithms, data, and decision-making processes
				\item Explaining specific predictions
				\item Reducing model size
			\end{enumerate}
		\end{block}
		\begin{exampleblock}{Answer: B}\end{exampleblock}
		\begin{block}{Explanation}
			Transparency = openness about design, algorithms, data, and processes. Distinct from interpretability (internal logic) and explainability (specific decisions).
		\end{block}
	\end{frame}
	
	\begin{frame}{Q69 --- Opaqueness Definition}
		\begin{block}{Question}
			Opaqueness in AI refers to:
			\begin{enumerate}
				\item Lack of transparency in a model's design, data, or decision-making
				\item Model generalisability
				\item Model accuracy
				\item Model speed
			\end{enumerate}
		\end{block}
		\begin{exampleblock}{Answer: A}\end{exampleblock}
		\begin{block}{Explanation}
			Opaqueness = lack of transparency. Forms include secrecy (unaware of profiling), confidentiality (internals hidden), and specialised knowledge (requires expertise).
		\end{block}
	\end{frame}
	
	\begin{frame}{Q70 --- Automation Bias Definition}
		\begin{block}{Question}
			Automation bias is best defined as:
			\begin{enumerate}
				\item A tendency to over-rely on automated systems, reducing vigilance
				\item Systematic measurement error
				\item Random sampling error
				\item Non-linear relationship in data
			\end{enumerate}
		\end{block}
		\begin{exampleblock}{Answer: A}\end{exampleblock}
		\begin{block}{Explanation}
			Automation bias = over-reliance on automated recommendations, reducing critical oversight and human verification.
		\end{block}
	\end{frame}
	
	\begin{frame}{Q71 --- Proxy Discrimination Definition}
		\begin{block}{Question}
			Proxy discrimination occurs when:
			\begin{enumerate}
				\item A feature correlated with a protected attribute carries discriminatory effects even after the protected attribute is removed
				\item The model is too complex
				\item The dataset is small
				\item The model uses too many features
			\end{enumerate}
		\end{block}
		\begin{exampleblock}{Answer: A}\end{exampleblock}
		\begin{block}{Explanation}
			Proxy discrimination = removing a protected attribute but retaining correlated features that perpetuate discrimination (e.g., ZIP codes as race proxies).
		\end{block}
	\end{frame}
	
	\begin{frame}{Q72 --- Historical Bias Definition}
		\begin{block}{Question}
			Historical bias in AI refers to:
			\begin{enumerate}
				\item Random noise in the training set
				\item Past discrimination encoded in training data, which the model then perpetuates
				\item Measurement error
				\item Sampling noise
			\end{enumerate}
		\end{block}
		\begin{exampleblock}{Answer: B}\end{exampleblock}
		\begin{block}{Explanation}
			Historical bias arises when training data reflects past discriminatory patterns. Models trained on such data reinforce and amplify these patterns.
		\end{block}
	\end{frame}
	
	\begin{frame}{Q73 --- Sampling Bias Definition}
		\begin{block}{Question}
			Sampling bias occurs when:
			\begin{enumerate}
				\item The sample is not representative of the target population
				\item The model is too simple
				\item The data has outliers
				\item The model overfits
			\end{enumerate}
		\end{block}
		\begin{exampleblock}{Answer: A}\end{exampleblock}
		\begin{block}{Explanation}
			Sampling bias = unrepresentative training sample. The model performs poorly when deployed on the actual target population.
		\end{block}
	\end{frame}
	
	\begin{frame}{Q74 --- Data Composition Bias Definition}
		\begin{block}{Question}
			Data composition bias arises from:
			\begin{enumerate}
				\item Imbalanced class distribution in the training data
				\item Features measured differently across groups
				\item Outliers in the data
				\item Model overfitting
			\end{enumerate}
		\end{block}
		\begin{exampleblock}{Answer: A}\end{exampleblock}
		\begin{block}{Explanation}
			Data composition bias = imbalance where minority classes are underrepresented. The model underperforms on minority groups.
		\end{block}
	\end{frame}
	
	\begin{frame}{Q75 --- Measurement Bias Definition}
		\begin{block}{Question}
			Measurement bias occurs when:
			\begin{enumerate}
				\item Features are measured differently across groups
				\item The class distribution is imbalanced
				\item The sample is small
				\item The model is nonlinear
			\end{enumerate}
		\end{block}
		\begin{exampleblock}{Answer: A}\end{exampleblock}
		\begin{block}{Explanation}
			Measurement bias = inconsistent measurement methods across groups (e.g., different diagnostic tests for different populations).
		\end{block}
	\end{frame}
	
	\begin{frame}{Q76 --- Individual Fairness Definition}
		\begin{block}{Question}
			Individual fairness is best defined as:
			\begin{enumerate}
				\item Equal outcomes across demographic groups
				\item Treating similar individuals similarly
				\item Equal true positive rates
				\item Equal precision
			\end{enumerate}
		\end{block}
		\begin{exampleblock}{Answer: B}\end{exampleblock}
		\begin{block}{Explanation}
			Individual fairness = treat like cases alike. Rooted in Aristotle's doctrine. Contrasts with group fairness (statistical parity).
		\end{block}
	\end{frame}
	
	\begin{frame}{Q77 --- Demographic Parity Definition}
		\begin{block}{Question}
			Demographic parity requires:
			\begin{enumerate}
				\item Equal true positive rates across groups
				\item Equal predicted positive rates across groups
				\item Equal precision across groups
				\item Equal accuracy across groups
			\end{enumerate}
		\end{block}
		\begin{exampleblock}{Answer: B}\end{exampleblock}
		\begin{block}{Explanation}
			Demographic parity: $P(\hat{Y}=1 | A=0) = P(\hat{Y}=1 | A=1)$. It does not consider whether predictions are correct.
		\end{block}
	\end{frame}
	
	\begin{frame}{Q78 --- Equal Opportunity Definition}
		\begin{block}{Question}
			Equal opportunity requires:
			\begin{enumerate}
				\item Equal positive prediction rates
				\item Equal true positive rates across groups
				\item Equal precision across groups
				\item Equal false positive rates
			\end{enumerate}
		\end{block}
		\begin{exampleblock}{Answer: B}\end{exampleblock}
		\begin{block}{Explanation}
			Equal opportunity: $P(\hat{Y}=1 | Y=1, A=0) = P(\hat{Y}=1 | Y=1, A=1)$. Requires equal TPR across groups.
		\end{block}
	\end{frame}
	
	\begin{frame}{Q79 --- Predictive Rate Parity Definition}
		\begin{block}{Question}
			Predictive rate parity requires:
			\begin{enumerate}
				\item Equal true positive rates across groups
				\item Equal precision across groups
				\item Equal positive prediction rates
				\item Equal accuracy
			\end{enumerate}
		\end{block}
		\begin{exampleblock}{Answer: B}\end{exampleblock}
		\begin{block}{Explanation}
			Predictive rate parity: $P(Y=1 | \hat{Y}=1, A=0) = P(Y=1 | \hat{Y}=1, A=1)$. Requires equal precision across groups.
		\end{block}
	\end{frame}
	
	\begin{frame}{Q80 --- Fairness Impossibility Theorem}
		\begin{block}{Question}
			The fairness impossibility theorem states that when base rates differ:
			\begin{enumerate}
				\item All fairness metrics can be satisfied
				\item Demographic parity, predictive rate parity, and equal opportunity cannot all be satisfied simultaneously
				\item Only equal opportunity can be satisfied
				\item Accuracy is the only relevant metric
			\end{enumerate}
		\end{block}
		\begin{exampleblock}{Answer: B}\end{exampleblock}
		\begin{block}{Explanation}
			Kleinberg, Mullainathan, and Raghavan (2016): under unequal base rates, these three fairness criteria are mutually incompatible.
		\end{block}
	\end{frame}
	
	\begin{frame}{Q81 --- Nonmaleficence Definition}
		\begin{block}{Question}
			Nonmaleficence refers to:
			\begin{enumerate}
				\item Actively promoting welfare
				\item Avoiding harm and unjustified risk
				\item Fairness
				\item Respecting autonomy
			\end{enumerate}
		\end{block}
		\begin{exampleblock}{Answer: B}\end{exampleblock}
		\begin{block}{Explanation}
			Nonmaleficence = do no harm. It forbids imposing unjustified risks.
		\end{block}
	\end{frame}
	
	\begin{frame}{Q82 --- Beneficence Definition}
		\begin{block}{Question}
			Beneficence refers to:
			\begin{enumerate}
				\item Avoiding harm
				\item Actively promoting the welfare of others
				\item Respecting autonomy
				\item Ensuring fairness
			\end{enumerate}
		\end{block}
		\begin{exampleblock}{Answer: B}\end{exampleblock}
		\begin{block}{Explanation}
			Beneficence = do good. It requires positive action to help others.
		\end{block}
	\end{frame}
	
	\begin{frame}{Q83 --- Autonomy Definition}
		\begin{block}{Question}
			Autonomy as an AI ethics principle requires:
			\begin{enumerate}
				\item Respecting individual decision-making
				\item Ensuring fairness
				\item Avoiding harm
				\item Benefiting others
			\end{enumerate}
		\end{block}
		\begin{exampleblock}{Answer: A}\end{exampleblock}
		\begin{block}{Explanation}
			Autonomy = respect for individual self-determination and consent.
		\end{block}
	\end{frame}
	
	\begin{frame}{Q84 --- Justice Definition}
		\begin{block}{Question}
			Justice as an AI ethics principle primarily concerns:
			\begin{enumerate}
				\item Fairness, equality, and impartiality
				\item Avoiding harm
				\item Benefiting others
				\item Respecting autonomy
			\end{enumerate}
		\end{block}
		\begin{exampleblock}{Answer: A}\end{exampleblock}
		\begin{block}{Explanation}
			Justice = fairness, equality, impartiality. Distributive justice concerns allocation; procedural justice concerns fair process.
		\end{block}
	\end{frame}
	
	\begin{frame}{Q85 --- Consequentialism Definition}
		\begin{block}{Question}
			Consequentialism is an ethical framework that judges actions by:
			\begin{enumerate}
				\item Adherence to rules
				\item Outcomes or consequences
				\item Character traits
				\item Legal compliance
			\end{enumerate}
		\end{block}
		\begin{exampleblock}{Answer: B}\end{exampleblock}
		\begin{block}{Explanation}
			Consequentialism judges by outcomes. Utilitarianism is the most common form, maximising overall welfare.
		\end{block}
	\end{frame}
	
	\begin{frame}{Q86 --- Deontology Definition}
		\begin{block}{Question}
			Deontology is an ethical framework that judges actions by:
			\begin{enumerate}
				\item Outcomes
				\item Adherence to duties and rules
				\item Character
				\item Profit
			\end{enumerate}
		\end{block}
		\begin{exampleblock}{Answer: B}\end{exampleblock}
		\begin{block}{Explanation}
			Deontology judges by duties and rules irrespective of consequences. Kant is the primary proponent.
		\end{block}
	\end{frame}
	
	\begin{frame}{Q87 --- Virtue Ethics Definition}
		\begin{block}{Question}
			Virtue ethics focuses on:
			\begin{enumerate}
				\item Outcomes
				\item Rules
				\item Character traits and living a good life
				\item Legal compliance
			\end{enumerate}
		\end{block}
		\begin{exampleblock}{Answer: C}\end{exampleblock}
		\begin{block}{Explanation}
			Virtue ethics asks what kind of person one should be and emphasises virtues like wisdom, courage, and justice.
		\end{block}
	\end{frame}
	
	\begin{frame}{Q88 --- Contextual Integrity Definition}
		\begin{block}{Question}
			Contextual integrity requires that:
			\begin{enumerate}
				\item Data can be freely reused
				\item Data use must respect the norms of the context in which it was provided
				\item Data must be encrypted
				\item Data must be deleted
			\end{enumerate}
		\end{block}
		\begin{exampleblock}{Answer: B}\end{exampleblock}
		\begin{block}{Explanation}
			Contextual integrity (Nissenbaum): data provided in one context should not be used in an incompatible context.
		\end{block}
	\end{frame}
	
	\begin{frame}{Q89 --- Data Minimisation Definition}
		\begin{block}{Question}
			Data minimisation requires:
			\begin{enumerate}
				\item Collecting all available data
				\item Collecting only data necessary for the specified purpose
				\item Sharing data freely
				\item Retaining data indefinitely
			\end{enumerate}
		\end{block}
		\begin{exampleblock}{Answer: B}\end{exampleblock}
		\begin{block}{Explanation}
			Data minimisation = collect only what is necessary for the specified purpose. Core GDPR principle.
		\end{block}
	\end{frame}
	
	\begin{frame}{Q90 --- Right to be Forgotten Definition}
		\begin{block}{Question}
			The right to be forgotten allows data subjects to:
			\begin{enumerate}
				\item Access their data
				\item Request erasure of personal data under certain conditions
				\item Sell their data
				\item Share data without consent
			\end{enumerate}
		\end{block}
		\begin{exampleblock}{Answer: B}\end{exampleblock}
		\begin{block}{Explanation}
			Right to erasure (GDPR Art. 17): allows deletion when data is no longer necessary, consent is withdrawn, etc.
		\end{block}
	\end{frame}
	
	\begin{frame}{Q91 --- Data Provenance Definition}
		\begin{block}{Question}
			Data provenance refers to:
			\begin{enumerate}
				\item Data storage format
				\item The origin, history, and legal right to use data
				\item Data encryption
				\item Data classification
			\end{enumerate}
		\end{block}
		\begin{exampleblock}{Answer: B}\end{exampleblock}
		\begin{block}{Explanation}
			Data provenance = origin and legal rights to use data. Critical for alternative data and AI training sets.
		\end{block}
	\end{frame}
	
	\begin{frame}{Q92 --- Data Classification Definition}
		\begin{block}{Question}
			Data classification is the process of:
			\begin{enumerate}
				\item Sorting data alphabetically
				\item Categorising data by sensitivity and structure to determine handling and controls
				\item Encrypting data
				\item Compressing data
			\end{enumerate}
		\end{block}
		\begin{exampleblock}{Answer: B}\end{exampleblock}
		\begin{block}{Explanation}
			Data classification assigns sensitivity labels (public, internal, confidential, restricted) and structure labels.
		\end{block}
	\end{frame}
	
	\begin{frame}{Q93 --- Metadata Definition}
		\begin{block}{Question}
			Metadata is best defined as:
			\begin{enumerate}
				\item Raw data
				\item Data about data
				\item Encrypted data
				\item Synthetic data
			\end{enumerate}
		\end{block}
		\begin{exampleblock}{Answer: B}\end{exampleblock}
		\begin{block}{Explanation}
			Metadata = data about data. It documents origin, structure, definitions, and relationships.
		\end{block}
	\end{frame}
	
	\begin{frame}{Q94 --- RBAC Definition}
		\begin{block}{Question}
			Role-Based Access Control (RBAC) grants access based on:
			\begin{enumerate}
				\item Individual discretion
				\item User roles within the organisation
				\item Random selection
				\item Data size
			\end{enumerate}
		\end{block}
		\begin{exampleblock}{Answer: B}\end{exampleblock}
		\begin{block}{Explanation}
			RBAC = permissions based on organisational roles. ABAC = attribute-based; MAC = mandatory; DAC = discretionary.
		\end{block}
	\end{frame}
	
	\begin{frame}{Q95 --- ABAC Definition}
		\begin{block}{Question}
			Attribute-Based Access Control (ABAC) grants access based on:
			\begin{enumerate}
				\item User roles
				\item Attributes of user, resource, and environment
				\item Random selection
				\item Seniority
			\end{enumerate}
		\end{block}
		\begin{exampleblock}{Answer: B}\end{exampleblock}
		\begin{block}{Explanation}
			ABAC uses attributes (user, resource, environment) to make access decisions. More flexible than RBAC.
		\end{block}
	\end{frame}
	
	\begin{frame}{Q96 --- Model Inventory Definition}
		\begin{block}{Question}
			A model inventory is:
			\begin{enumerate}
				\item A list of all models in use, with owners, risk tiers, and validation status
				\item A list of data sources
				\item A list of employees
				\item A list of IT systems
			\end{enumerate}
		\end{block}
		\begin{exampleblock}{Answer: A}\end{exampleblock}
		\begin{block}{Explanation}
			Model inventory = registry of all models. Includes owner, purpose, risk tier, validation status, and interdependencies.
		\end{block}
	\end{frame}
	
	\begin{frame}{Q97 --- Model Landscape Definition}
		\begin{block}{Question}
			A model landscape is:
			\begin{enumerate}
				\item A list of all models
				\item A map of interactions and interdependencies between models
				\item A training dataset
				\item A model architecture
			\end{enumerate}
		\end{block}
		\begin{exampleblock}{Answer: B}\end{exampleblock}
		\begin{block}{Explanation}
			Model landscape = map of interactions and dependencies. It reveals cascade risk and critical nodes.
		\end{block}
	\end{frame}
	
	\begin{frame}{Q98 --- Model Risk Appetite Definition}
		\begin{block}{Question}
			Model risk appetite is:
			\begin{enumerate}
				\item The set of models in production
				\item The level of model risk the organisation is willing to accept
				\item The cost of modelling
				\item The number of models
			\end{enumerate}
		\end{block}
		\begin{exampleblock}{Answer: B}\end{exampleblock}
		\begin{block}{Explanation}
			Model risk appetite = acceptable level of model risk, set by the board, guiding governance.
		\end{block}
	\end{frame}
	
	\begin{frame}{Q99 --- Use Test Definition}
		\begin{block}{Question}
			The use test in model validation is:
			\begin{enumerate}
				\item A quantitative accuracy metric
				\item A qualitative assessment of whether the model is used appropriately for its intended purpose
				\item A code review
				\item A speed test
			\end{enumerate}
		\end{block}
		\begin{exampleblock}{Answer: B}\end{exampleblock}
		\begin{block}{Explanation}
			The use test assesses whether the model is used correctly and appropriately in practice. It is qualitative and people-focused.
		\end{block}
	\end{frame}
	
	\begin{frame}{Q100 --- Outcomes Analysis Definition}
		\begin{block}{Question}
			Outcomes analysis in model validation is used to:
			\begin{enumerate}
				\item Speed up training
				\item Compare model predictions against realised outcomes to detect bias, degradation, or incorrect assumptions
				\item Reduce features
				\item Increase accuracy
			\end{enumerate}
		\end{block}
		\begin{exampleblock}{Answer: B}\end{exampleblock}
		\begin{block}{Explanation}
			Outcomes analysis compares model predictions against actual outcomes to validate model quality in production.
		\end{block}
	\end{frame}
	
	% =================================================================
	\section{Algorithm Names and Mechanism Recall}
	% =================================================================
	
	\begin{frame}{Q101 --- K-means++ Purpose}
		\begin{block}{Question}
			K-means++ is an improved initialisation method whose purpose is to:
			\begin{enumerate}
				\item Reduce K
				\item Spread initial centroids far apart to reduce poor local optima
				\item Increase WCSS
				\item Eliminate outliers
			\end{enumerate}
		\end{block}
		\begin{exampleblock}{Answer: B}\end{exampleblock}
		\begin{block}{Explanation}
			K-means++ selects initial centroids with probability proportional to squared distance from nearest existing centroid, spreading them out and reducing poor local optima.
		\end{block}
	\end{frame}
	
	\begin{frame}{Q102 --- Scree Plot Definition}
		\begin{block}{Question}
			A scree plot displays:
			\begin{enumerate}
				\item WCSS against number of clusters K
				\item Data against time
				\item Two features against each other
				\item Model accuracy against training time
			\end{enumerate}
		\end{block}
		\begin{exampleblock}{Answer: A}\end{exampleblock}
		\begin{block}{Explanation}
			A scree plot charts WCSS (inertia) against K. The elbow point suggests the optimal K.
		\end{block}
	\end{frame}
	
	\begin{frame}{Q103 --- Dendrogram Definition}
		\begin{block}{Question}
			A dendrogram is:
			\begin{enumerate}
				\item A scatter plot
				\item A tree-like diagram showing hierarchical clustering steps
				\item A line chart of loss
				\item A histogram
			\end{enumerate}
		\end{block}
		\begin{exampleblock}{Answer: B}\end{exampleblock}
		\begin{block}{Explanation}
			A dendrogram visualises hierarchical clustering: x-axis shows data points, y-axis shows merge distance. Long vertical lines indicate significant cluster boundaries.
		\end{block}
	\end{frame}
	
	\begin{frame}{Q104 --- Confusion Matrix Layout}
		\begin{block}{Question}
			A standard confusion matrix contains four cells:
			\begin{enumerate}
				\item TP, TN, FP, FN
				\item A, B, C, D
				\item Alpha, Beta, Gamma, Delta
				\item Mean, Median, Mode, Range
			\end{enumerate}
		\end{block}
		\begin{exampleblock}{Answer: A}\end{exampleblock}
		\begin{block}{Explanation}
			Confusion matrix cells: True Positive, True Negative, False Positive (Type I error), False Negative (Type II error).
		\end{block}
	\end{frame}
	
	\begin{frame}{Q105 --- ROC Axes}
		\begin{block}{Question}
			The ROC curve plots which two quantities on its axes?
			\begin{enumerate}
				\item Precision vs recall
				\item True positive rate vs false positive rate
				\item Train error vs test error
				\item Loss vs iteration
			\end{enumerate}
		\end{block}
		\begin{exampleblock}{Answer: B}\end{exampleblock}
		\begin{block}{Explanation}
			ROC: y-axis = TPR (sensitivity); x-axis = FPR (1 $-$ specificity). AUC summarises discrimination.
		\end{block}
	\end{frame}
	
	\begin{frame}{Q106 --- AUC Range and Interpretation}
		\begin{block}{Question}
			The AUC ranges over which values, and what does AUC = 0.5 indicate?
			\begin{enumerate}
				\item $[0, 1]$; AUC = 0.5 indicates random classifier
				\item $[-1, 1]$; AUC = 0.5 indicates perfect classifier
				\item $[0, 100]$; AUC = 0.5 indicates perfect classifier
				\item $[1, 10]$; AUC = 0.5 indicates random classifier
			\end{enumerate}
		\end{block}
		\begin{exampleblock}{Answer: A}\end{exampleblock}
		\begin{block}{Explanation}
			AUC ranges in $[0, 1]$. AUC = 0.5 = random; AUC = 1.0 = perfect; AUC $<$ 0.5 = worse than random.
		\end{block}
	\end{frame}
	
	\begin{frame}{Q107 --- K-Fold Cross-Validation Definition}
		\begin{block}{Question}
			In $k$-fold cross-validation, the data is:
			\begin{enumerate}
				\item Split into $k$ equal folds; the model is trained $k$ times, each time holding out a different fold for validation
				\item Split into 2 folds only
				\item Not split at all
				\item Split randomly without reuse
			\end{enumerate}
		\end{block}
		\begin{exampleblock}{Answer: A}\end{exampleblock}
		\begin{block}{Explanation}
			$k$-fold CV: partition into $k$ folds; iterate $k$ times, each time training on $k-1$ folds and validating on the held-out fold. Average the results.
		\end{block}
	\end{frame}
	
	\begin{frame}{Q108 --- Stratified CV Purpose}
		\begin{block}{Question}
			Stratified cross-validation ensures that:
			\begin{enumerate}
				\item Each fold has the same class proportion as the whole dataset
				\item Only the majority class is used
				\item All folds have equal mean
				\item Random assignment is avoided
			\end{enumerate}
		\end{block}
		\begin{exampleblock}{Answer: A}\end{exampleblock}
		\begin{block}{Explanation}
			Stratified CV preserves class ratios in every fold, critical for imbalanced datasets.
		\end{block}
	\end{frame}
	
	\begin{frame}{Q109 --- Bootstrap Definition}
		\begin{block}{Question}
			Bootstrapping in model evaluation is:
			\begin{enumerate}
				\item Sampling with replacement to create new datasets
				\item Sampling without replacement
				\item Splitting data into two halves
				\item Training without validation
			\end{enumerate}
		\end{block}
		\begin{exampleblock}{Answer: A}\end{exampleblock}
		\begin{block}{Explanation}
			Bootstrapping draws samples with replacement. About 63\% unique observations per sample; ~37\% out-of-bootstrap.
		\end{block}
	\end{frame}
	
	\begin{frame}{Q110 --- Out-of-Bag Fraction}
		\begin{block}{Question}
			In bootstrapping, the expected fraction of observations NOT sampled in each bootstrap sample is approximately:
			\begin{enumerate}
				\item 10\%
				\item 25\%
				\item 37\%
				\item 63\%
			\end{enumerate}
		\end{block}
		\begin{exampleblock}{Answer: C}\end{exampleblock}
		\begin{block}{Explanation}
			The expected out-of-bootstrap fraction is $1 - 1/e \approx 0.368$ = 37\%.
		\end{block}
	\end{frame}
	
	\begin{frame}{Q111 --- Grid Search Definition}
		\begin{block}{Question}
			Grid search is a technique for:
			\begin{enumerate}
				\item Exhaustively searching through a predefined set of hyperparameter combinations
				\item Removing outliers
				\item Feature selection only
				\item Clustering
			\end{enumerate}
		\end{block}
		\begin{exampleblock}{Answer: A}\end{exampleblock}
		\begin{block}{Explanation}
			Grid search evaluates every combination of hyperparameters within specified ranges. Often combined with cross-validation.
		\end{block}
	\end{frame}
	
	\begin{frame}{Q112 --- Learning Rate Definition}
		\begin{block}{Question}
			The learning rate $\eta$ in gradient descent controls:
			\begin{enumerate}
				\item The number of iterations
				\item The step size of each weight update
				\item The number of features
				\item The batch size
			\end{enumerate}
		\end{block}
		\begin{exampleblock}{Answer: B}\end{exampleblock}
		\begin{block}{Explanation}
			Learning rate $\eta$ scales the gradient to determine step size. Too small = slow; too large = overshoot.
		\end{block}
	\end{frame}
	
	\begin{frame}{Q113 --- Backpropagation Definition}
		\begin{block}{Question}
			Backpropagation is:
			\begin{enumerate}
				\item Forward-only computation
				\item An algorithm that computes gradients layer-by-layer using the chain rule, starting from the output layer
				\item A clustering method
				\item A type of regularisation
			\end{enumerate}
		\end{block}
		\begin{exampleblock}{Answer: B}\end{exampleblock}
		\begin{block}{Explanation}
			Backpropagation computes gradients efficiently by applying the chain rule from output to input layers.
		\end{block}
	\end{frame}
	
	\begin{frame}{Q114 --- Epoch Definition}
		\begin{block}{Question}
			An epoch in neural network training is:
			\begin{enumerate}
				\item One gradient update
				\item One full pass through the entire training dataset
				\item One validation step
				\item One feature selection step
			\end{enumerate}
		\end{block}
		\begin{exampleblock}{Answer: B}\end{exampleblock}
		\begin{block}{Explanation}
			An epoch = one full pass through the training data. Typically many epochs are required for convergence.
		\end{block}
	\end{frame}
	
	\begin{frame}{Q115 --- Batch Gradient Descent Definition}
		\begin{block}{Question}
			Batch gradient descent updates weights using:
			\begin{enumerate}
				\item One random training example per step
				\item A subset of training examples per step
				\item The entire training set per step
				\item No training examples
			\end{enumerate}
		\end{block}
		\begin{exampleblock}{Answer: C}\end{exampleblock}
		\begin{block}{Explanation}
			Batch GD uses the full dataset per update. Stochastic uses one example; mini-batch uses a subset.
		\end{block}
	\end{frame}
	
	\begin{frame}{Q116 --- Stochastic Gradient Descent Definition}
		\begin{block}{Question}
			Stochastic gradient descent (SGD) updates weights using:
			\begin{enumerate}
				\item The full training set
				\item One random training example per step
				\item A fixed subset
				\item No examples
			\end{enumerate}
		\end{block}
		\begin{exampleblock}{Answer: B}\end{exampleblock}
		\begin{block}{Explanation}
			SGD uses one randomly selected example per update. It is noisier but faster per pass.
		\end{block}
	\end{frame}
	
	\begin{frame}{Q117 --- Mini-Batch Gradient Descent Definition}
		\begin{block}{Question}
			Mini-batch gradient descent updates weights using:
			\begin{enumerate}
				\item One example
				\item The full dataset
				\item A small subset of the training data per step
				\item No data
			\end{enumerate}
		\end{block}
		\begin{exampleblock}{Answer: C}\end{exampleblock}
		\begin{block}{Explanation}
			Mini-batch GD uses a small subset (e.g., 32 or 64 examples). It balances parallelism and stochasticity.
		\end{block}
	\end{frame}
	
	\begin{frame}{Q118 --- Momentum Purpose}
		\begin{block}{Question}
			Momentum in gradient descent primarily:
			\begin{enumerate}
				\item Reduces the learning rate
				\item Smooths the optimisation path and accelerates in consistent gradient directions
				\item Eliminates local minima
				\item Converts non-convex to convex
			\end{enumerate}
		\end{block}
		\begin{exampleblock}{Answer: B}\end{exampleblock}
		\begin{block}{Explanation}
			Momentum accumulates a velocity vector, accelerating consistent directions and dampening oscillation.
		\end{block}
	\end{frame}
	
	\begin{frame}{Q119 --- Adam Optimiser}
		\begin{block}{Question}
			The Adam optimiser combines which two ideas?
			\begin{enumerate}
				\item Momentum and adaptive learning rates
				\item Batch and stochastic descent
				\item Ridge and LASSO
				\item Precision and recall
			\end{enumerate}
		\end{block}
		\begin{exampleblock}{Answer: A}\end{exampleblock}
		\begin{block}{Explanation}
			Adam = Adaptive Moment Estimation. It combines momentum (first moment) with adaptive per-parameter learning rates (second moment).
		\end{block}
	\end{frame}
	
	\begin{frame}{Q120 --- Dropout Purpose}
		\begin{block}{Question}
			Dropout regularisation works by:
			\begin{enumerate}
				\item Removing whole layers
				\item Randomly deactivating a fraction of neurons during training to prevent co-adaptation
				\item Reducing the learning rate
				\item Removing features
			\end{enumerate}
		\end{block}
		\begin{exampleblock}{Answer: B}\end{exampleblock}
		\begin{block}{Explanation}
			Dropout randomly drops units during training. At inference, all units are used with scaled weights.
		\end{block}
	\end{frame}
	
	\begin{frame}{Q121 --- Batch Normalization Purpose}
		\begin{block}{Question}
			Batch normalisation normalises:
			\begin{enumerate}
				\item Weights
				\item Intermediate layer inputs, stabilising training and mitigating internal covariate shift
				\item Features
				\item Outputs
			\end{enumerate}
		\end{block}
		\begin{exampleblock}{Answer: B}\end{exampleblock}
		\begin{block}{Explanation}
			Batch norm normalises the activations of a layer across the batch, improving stability and speed of training.
		\end{block}
	\end{frame}
	
	\begin{frame}{Q122 --- Early Stopping Purpose}
		\begin{block}{Question}
			Early stopping is a regularisation technique that:
			\begin{enumerate}
				\item Stops training when validation performance starts to deteriorate
				\item Stops training after a fixed number of epochs
				\item Stops training immediately
				\item Never stops training
			\end{enumerate}
		\end{block}
		\begin{exampleblock}{Answer: A}\end{exampleblock}
		\begin{block}{Explanation}
			Early stopping monitors validation loss and halts training when it stops improving, preventing overfitting.
		\end{block}
	\end{frame}
	
	\begin{frame}{Q123 --- Vanishing Gradient Definition}
		\begin{block}{Question}
			The vanishing gradient problem occurs when:
			\begin{enumerate}
				\item Gradients grow unboundedly
				\item Gradients become extremely small as they propagate through many layers
				\item Weights are initialised poorly
				\item The learning rate is too high
			\end{enumerate}
		\end{block}
		\begin{exampleblock}{Answer: B}\end{exampleblock}
		\begin{block}{Explanation}
			Vanishing gradient: gradients shrink exponentially through deep networks with saturating activations (e.g., sigmoid).
		\end{block}
	\end{frame}
	
	\begin{frame}{Q124 --- Exploding Gradient Definition}
		\begin{block}{Question}
			The exploding gradient problem occurs when:
			\begin{enumerate}
				\item Gradients shrink to near zero
				\item Gradients grow extremely large across layers, destabilising training
				\item Weights are zero
				\item Learning rate is zero
			\end{enumerate}
		\end{block}
		\begin{exampleblock}{Answer: B}\end{exampleblock}
		\begin{block}{Explanation}
			Exploding gradient: gradients grow exponentially. Remedies include gradient clipping and careful initialisation.
		\end{block}
	\end{frame}
	
	\begin{frame}{Q125 --- ReLU Activation}
		\begin{block}{Question}
			The ReLU (Rectified Linear Unit) activation function is:
			\begin{enumerate}
				\item $f(x) = \max(0, x)$
				\item $f(x) = 1/(1+e^{-x})$
				\item $f(x) = \tanh(x)$
				\item $f(x) = x^2$
			\end{enumerate}
		\end{block}
		\begin{exampleblock}{Answer: A}\end{exampleblock}
		\begin{block}{Explanation}
			ReLU = $\max(0, x)$. It is computationally efficient and reduces vanishing gradient for positive values.
		\end{block}
	\end{frame}
	
	\begin{frame}{Q126 --- Sigmoid Activation}
		\begin{block}{Question}
			The sigmoid activation function is:
			\begin{enumerate}
				\item $f(x) = \max(0, x)$
				\item $f(x) = 1/(1+e^{-x})$
				\item $f(x) = \tanh(x)$
				\item $f(x) = x$
			\end{enumerate}
		\end{block}
		\begin{exampleblock}{Answer: B}\end{exampleblock}
		\begin{block}{Explanation}
			Sigmoid = $1/(1+e^{-x})$, bounded in $(0, 1)$. Used in binary outputs but can suffer from vanishing gradients.
		\end{block}
	\end{frame}
	
	\begin{frame}{Q127 --- Tanh Activation}
		\begin{block}{Question}
			The tanh (hyperbolic tangent) activation function is:
			\begin{enumerate}
				\item $f(x) = \max(0, x)$
				\item $f(x) = 1/(1+e^{-x})$
				\item $f(x) = (e^x - e^{-x})/(e^x + e^{-x})$
				\item $f(x) = x^2$
			\end{enumerate}
		\end{block}
		\begin{exampleblock}{Answer: C}\end{exampleblock}
		\begin{block}{Explanation}
			tanh = $(e^x - e^{-x})/(e^x + e^{-x})$. Output is bounded in $[-1, 1]$ and centred at zero.
		\end{block}
	\end{frame}
	
	\begin{frame}{Q128 --- Softmax Purpose}
		\begin{block}{Question}
			Softmax is used at the output layer primarily for:
			\begin{enumerate}
				\item Binary classification
				\item Multi-class classification, converting logits into a probability distribution
				\item Regression
				\item Clustering
			\end{enumerate}
		\end{block}
		\begin{exampleblock}{Answer: B}\end{exampleblock}
		\begin{block}{Explanation}
			Softmax outputs sum to one, making it appropriate for multi-class probability outputs.
		\end{block}
	\end{frame}
	
	\begin{frame}{Q129 --- Confusion Matrix Type I Error}
		\begin{block}{Question}
			A Type I error in a classification context corresponds to:
			\begin{enumerate}
				\item False positive
				\item False negative
				\item True positive
				\item True negative
			\end{enumerate}
		\end{block}
		\begin{exampleblock}{Answer: A}\end{exampleblock}
		\begin{block}{Explanation}
			Type I error = false positive (rejecting a true null hypothesis). Type II error = false negative.
		\end{block}
	\end{frame}
	
	\begin{frame}{Q130 --- Confusion Matrix Type II Error}
		\begin{block}{Question}
			A Type II error in a classification context corresponds to:
			\begin{enumerate}
				\item False positive
				\item False negative
				\item True positive
				\item True negative
			\end{enumerate}
		\end{block}
		\begin{exampleblock}{Answer: B}\end{exampleblock}
		\begin{block}{Explanation}
			Type II error = false negative (failing to reject a false null hypothesis).
		\end{block}
	\end{frame}
	
	\begin{frame}{Q131 --- Class Imbalance Handling Techniques}
		\begin{block}{Question}
			Which of the following is NOT a standard technique for handling class imbalance?
			\begin{enumerate}
				\item Oversampling the minority class
				\item Undersampling the majority class
				\item Using a weighted loss function
				\item Removing the minority class entirely
			\end{enumerate}
		\end{block}
		\begin{exampleblock}{Answer: D}\end{exampleblock}
		\begin{block}{Explanation}
			Removing the minority class destroys the signal and defeats the purpose. Options A, B, C are standard approaches.
		\end{block}
	\end{frame}
	
	\begin{frame}{Q132 --- SMOTE Definition}
		\begin{block}{Question}
			SMOTE stands for:
			\begin{enumerate}
				\item Synthetic Minority Over-sampling Technique
				\item Statistical Missing Value Estimation
				\item Sparse Matrix Optimisation Technique
				\item Sequential Model Ensemble
			\end{enumerate}
		\end{block}
		\begin{exampleblock}{Answer: A}\end{exampleblock}
		\begin{block}{Explanation}
			SMOTE = Synthetic Minority Over-sampling Technique. It generates synthetic minority samples via interpolation between existing minority points.
		\end{block}
	\end{frame}
	
	\begin{frame}{Q133 --- Bagging Definition}
		\begin{block}{Question}
			Bagging (Bootstrap Aggregation) is:
			\begin{enumerate}
				\item Training models sequentially
				\item Training multiple models on bootstrap samples and averaging predictions
				\item Training a single model on all data
				\item Removing features
			\end{enumerate}
		\end{block}
		\begin{exampleblock}{Answer: B}\end{exampleblock}
		\begin{block}{Explanation}
			Bagging trains models in parallel on bootstrap samples and averages outputs. Reduces variance.
		\end{block}
	\end{frame}
	
	\begin{frame}{Q134 --- Random Forest Definition}
		\begin{block}{Question}
			Random forest differs from bagging by:
			\begin{enumerate}
				\item Using sequential training
				\item Adding random feature subsetting at each split
				\item Using only one tree
				\item Removing randomness
			\end{enumerate}
		\end{block}
		\begin{exampleblock}{Answer: B}\end{exampleblock}
		\begin{block}{Explanation}
			Random forest = bagging + random feature subsampling at each split. Reduces correlation across trees.
		\end{block}
	\end{frame}
	
	\begin{frame}{Q135 --- Boosting Definition}
		\begin{block}{Question}
			Boosting trains models:
			\begin{enumerate}
				\item In parallel on bootstrap samples
				\item Sequentially, each correcting the errors of the previous model
				\item Randomly
				\item Without any error correction
			\end{enumerate}
		\end{block}
		\begin{exampleblock}{Answer: B}\end{exampleblock}
		\begin{block}{Explanation}
			Boosting fits models sequentially, each new model focused on prior residuals or misclassified points. Reduces bias.
		\end{block}
	\end{frame}
	
	\begin{frame}{Q136 --- AdaBoost Definition}
		\begin{block}{Question}
			AdaBoost is a boosting method that:
			\begin{enumerate}
				\item Fits residuals
				\item Reweights misclassified training examples and combines weak learners
				\item Uses random features
				\item Removes features
			\end{enumerate}
		\end{block}
		\begin{exampleblock}{Answer: B}\end{exampleblock}
		\begin{block}{Explanation}
			AdaBoost reweights training examples after each iteration, focusing subsequent weak learners on previously misclassified points.
		\end{block}
	\end{frame}
	
	\begin{frame}{Q137 --- Gradient Boosting Definition}
		\begin{block}{Question}
			Gradient boosting fits new trees to:
			\begin{enumerate}
				\item The original labels
				\item The residuals of previous trees
				\item Random labels
				\item Boostrap samples only
			\end{enumerate}
		\end{block}
		\begin{exampleblock}{Answer: B}\end{exampleblock}
		\begin{block}{Explanation}
			Gradient boosting fits each new tree to the negative gradient (residuals) of the loss function from prior trees.
		\end{block}
	\end{frame}
	
	\begin{frame}{Q138 --- XGBoost Definition}
		\begin{block}{Question}
			XGBoost stands for:
			\begin{enumerate}
				\item Extreme Gradient Boosting
				\item Extended Gradient Boost
				\item Extra Gradient Boost
				\item Exponential Gradient Boost
			\end{enumerate}
		\end{block}
		\begin{exampleblock}{Answer: A}\end{exampleblock}
		\begin{block}{Explanation}
			XGBoost = Extreme Gradient Boosting. An efficient implementation of gradient boosting with regularisation.
		\end{block}
	\end{frame}
	
	\begin{frame}{Q139 --- Bagging vs Boosting Variance Bias}
		\begin{block}{Question}
			Which statement correctly distinguishes bagging and boosting in terms of bias and variance?
			\begin{enumerate}
				\item Bagging mainly reduces variance; boosting mainly reduces bias
				\item Bagging mainly reduces bias; boosting mainly reduces variance
				\item Both reduce bias only
				\item Both reduce variance only
			\end{enumerate}
		\end{block}
		\begin{exampleblock}{Answer: A}\end{exampleblock}
		\begin{block}{Explanation}
			Bagging averages decorrelated models, reducing variance. Boosting sequentially corrects errors, reducing bias.
		\end{block}
	\end{frame}
	
	\begin{frame}{Q140 --- Precision-Recall Trade-off}
		\begin{block}{Question}
			Lowering the classification threshold generally:
			\begin{enumerate}
				\item Increases recall and decreases precision
				\item Increases precision and decreases recall
				\item Increases both
				\item Decreases both
			\end{enumerate}
		\end{block}
		\begin{exampleblock}{Answer: A}\end{exampleblock}
		\begin{block}{Explanation}
			Lower threshold = more positives predicted = higher recall but lower precision (more false positives).
		\end{block}
	\end{frame}
	
	% =================================================================
	\section{Hard Theory and Formula Questions (continued)}
	% =================================================================
	
	\begin{frame}{Q141 --- Log-Odds in Logistic Regression}
		\begin{block}{Question}
			In logistic regression, the log-odds (logit) of the outcome is:
			\begin{enumerate}
				\item $\ln(P/(1-P)) = \beta_0 + \beta_1 x_1 + \cdots$
				\item $P/(1-P)$
				\item $P \cdot (1-P)$
				\item $\ln(P)$
			\end{enumerate}
		\end{block}
		\begin{exampleblock}{Answer: A}\end{exampleblock}
		\begin{block}{Explanation}
			Logit = $\ln(P/(1-P))$, linear in the coefficients. Logistic regression models the log-odds as a linear function of features.
		\end{block}
	\end{frame}
	
	\begin{frame}{Q142 --- Odds Ratio Interpretation}
		\begin{block}{Question}
			In logistic regression, the exponential of a coefficient $\exp(\beta_j)$ represents:
			\begin{enumerate}
				\item The probability
				\item The odds ratio for a one-unit increase in $x_j$
				\item The variance
				\item The residual
			\end{enumerate}
		\end{block}
		\begin{exampleblock}{Answer: B}\end{exampleblock}
		\begin{block}{Explanation}
			$\exp(\beta_j)$ = odds ratio. A one-unit increase in $x_j$ multiplies the odds by $\exp(\beta_j)$.
		\end{block}
	\end{frame}
	
	\begin{frame}{Q143 --- Linear Discriminant Analysis Formula}
		\begin{block}{Question}
			In LDA, the discriminant function for class $j$ is:
			\begin{enumerate}
				\item $\delta_j(x) = x^T \hat{\Sigma}^{-1} \hat{\mu}_j - \frac{1}{2}\hat{\mu}_j^T \hat{\Sigma}^{-1} \hat{\mu}_j + \ln(\hat{\pi}_j)$
				\item $\delta_j(x) = x$
				\item $\delta_j(x) = x^2$
				\item $\delta_j(x) = 1/(1+e^{-x})$
			\end{enumerate}
		\end{block}
		\begin{exampleblock}{Answer: A}\end{exampleblock}
		\begin{block}{Explanation}
			LDA discriminant: $\delta_j(x) = x^T \hat{\Sigma}^{-1}\hat{\mu}_j - \frac{1}{2}\hat{\mu}_j^T \hat{\Sigma}^{-1}\hat{\mu}_j + \ln(\hat{\pi}_j)$. Classification chooses the largest $\delta_j$.
		\end{block}
	\end{frame}
	
	\begin{frame}{Q144 --- Covariance Matrix Inverse (2x2)}
		\begin{block}{Question}
			For a 2x2 matrix $\begin{pmatrix} a & b \\ c & d \end{pmatrix}$, the inverse is:
			\begin{enumerate}
				\item $\dfrac{1}{ad - bc} \begin{pmatrix} d & -b \\ -c & a \end{pmatrix}$
				\item $\begin{pmatrix} d & -b \\ -c & a \end{pmatrix}$
				\item $\begin{pmatrix} a & b \\ c & d \end{pmatrix}$
				\item $\dfrac{1}{2} \begin{pmatrix} a & b \\ c & d \end{pmatrix}$
			\end{enumerate}
		\end{block}
		\begin{exampleblock}{Answer: A}\end{exampleblock}
		\begin{block}{Explanation}
			Inverse of 2x2 = 1/determinant $\times$ adjugate. Determinant = $ad - bc$.
		\end{block}
	\end{frame}
	
	\begin{frame}{Q145 --- Covariance Formula}
		\begin{block}{Question}
			The covariance between two features $x_1$ and $x_2$ is:
			\begin{enumerate}
				\item $\dfrac{\sum (x_1 - \bar{x}_1)(x_2 - \bar{x}_2)}{N - 1}$
				\item $\sum x_1 x_2$
				\item $\bar{x}_1 \cdot \bar{x}_2$
				\item $\dfrac{\sum x_1}{N} \cdot \dfrac{\sum x_2}{N}$
			\end{enumerate}
		\end{block}
		\begin{exampleblock}{Answer: A}\end{exampleblock}
		\begin{block}{Explanation}
			Covariance = average product of deviations from means, using $N-1$ for sample covariance.
		\end{block}
	\end{frame}
	
	\begin{frame}{Q146 --- AIC Formula}
		\begin{block}{Question}
			The Akaike Information Criterion is defined as:
			\begin{enumerate}
				\item $AIC = 2k - 2\ln(L)$
				\item $AIC = k + L$
				\item $AIC = 2k + 2\ln(L)$
				\item $AIC = L/k$
			\end{enumerate}
		\end{block}
		\begin{exampleblock}{Answer: A}\end{exampleblock}
		\begin{block}{Explanation}
			AIC = $2k - 2\ln(L)$ where $k$ = number of parameters and $L$ = likelihood. Lower AIC = better model.
		\end{block}
	\end{frame}
	
	\begin{frame}{Q147 --- BIC Formula}
		\begin{block}{Question}
			The Bayesian Information Criterion is defined as:
			\begin{enumerate}
				\item $BIC = k\ln(N) - 2\ln(L)$
				\item $BIC = 2k - 2\ln(L)$
				\item $BIC = k + L$
				\item $BIC = L^2$
			\end{enumerate}
		\end{block}
		\begin{exampleblock}{Answer: A}\end{exampleblock}
		\begin{block}{Explanation}
			BIC = $k\ln(N) - 2\ln(L)$. It penalises complexity more heavily than AIC for large samples.
		\end{block}
	\end{frame}
	
	\begin{frame}{Q148 --- Heteroskedasticity Definition}
		\begin{block}{Question}
			Heteroskedasticity in regression refers to:
			\begin{enumerate}
				\item Non-constant variance of residuals across observations
				\item Non-normal residuals
				\item Nonlinearity in relationships
				\item Multicollinearity
			\end{enumerate}
		\end{block}
		\begin{exampleblock}{Answer: A}\end{exampleblock}
		\begin{block}{Explanation}
			Heteroskedasticity = variance of errors is not constant. It inflates standard errors and invalidates inference.
		\end{block}
	\end{frame}
	
	\begin{frame}{Q149 --- Cook's Distance Purpose}
		\begin{block}{Question}
			Cook's distance is used to:
			\begin{enumerate}
				\item Measure the influence of each observation on fitted regression coefficients
				\item Measure model accuracy
				\item Detect multicollinearity
				\item Estimate variance
			\end{enumerate}
		\end{block}
		\begin{exampleblock}{Answer: A}\end{exampleblock}
		\begin{block}{Explanation}
			Cook's distance measures how much coefficients change if an observation is removed. High values identify influential points.
		\end{block}
	\end{frame}
	
	\begin{frame}{Q150 --- White's Test Purpose}
		\begin{block}{Question}
			White's Test in regression is used to:
			\begin{enumerate}
				\item Test for heteroskedasticity
				\item Test for normality
				\item Test for autocorrelation
				\item Test for linearity
			\end{enumerate}
		\end{block}
		\begin{exampleblock}{Answer: A}\end{exampleblock}
		\begin{block}{Explanation}
			White's Test detects heteroskedasticity (non-constant error variance). A significant result indicates heteroskedasticity.
		\end{block}
	\end{frame}
	
\end{document}